\documentclass[11pt]{article}
\usepackage{amsmath, amssymb, amsthm}

\title{Worst-Case Ratio for $P||\sum p_j C_j(\alpha)$ under Pair-Separation Restriction}
\author{}
\date{}

\begin{document}

\maketitle

\section{Problem Definition}

We consider the scheduling problem on $m$ identical parallel machines. Let there be $n$ jobs, where each job $j$ has processing time $p_j > 0$.

For a given schedule $S$, define the modified completion time of job $j$ as
\[
C_j(\alpha) = B_j + \alpha p_j,
\]
where $B_j$ is the total processing time of jobs scheduled before $j$ on the same machine, and $\alpha \in (0,1]$ is a fixed parameter.

The objective function is
\[
F_\alpha(S) = \sum_{j=1}^n p_j C_j(\alpha).
\]

\section{Pairwise Reformulation}

For any schedule $S$, the objective can be rewritten as
\[
F_\alpha(S)
=
\alpha \sum_{j=1}^n p_j^2
+
\sum_{\{j,k\}: j,k \text{ on same machine in } S} p_j p_k.
\]

Define
\[
S_2 = \sum_{j=1}^n p_j^2.
\]

Thus,
\[
F_\alpha(S) = \alpha S_2 + \text{pair contribution}.
\]

\section{Worst Schedule and Restricted Optimal Schedule}

Let $W$ denote a worst-case schedule maximizing $F_\alpha(S)$.

Let $OPT^R$ denote an optimal schedule subject to the following restriction:

\medskip

\noindent
\textbf{Pair-Separation Restriction:}
\begin{quote}
If two jobs $j,k$ are scheduled on the same machine in $W$, then they are not allowed to be scheduled on the same machine in $OPT^R$.
\end{quote}

\medskip

\section{Structural Consequence}

Under this restriction, any set of jobs that are assigned to the same machine in $W$ must be fully separated across machines in $OPT^R$.

Since $OPT^R$ uses only $m$ machines, it follows that no machine in $W$ can contain more than $m$ jobs. Otherwise, if a machine in $W$ contained $m+1$ jobs, these jobs would require at least $m+1$ machines to be pairwise separated in $OPT^R$, which is impossible.

Therefore, for every machine $i$ in $W$,
\[
|S_i^W| \le m.
\]

\section{Bounding the Worst Schedule}

Fix a machine $i$ in $W$ with job set $S_i^W$. Let
\[
L_i = \sum_{j \in S_i^W} p_j, \qquad r_i = |S_i^W| \le m.
\]

The pair contribution on machine $i$ is
\[
\sum_{\{j,k\} \subseteq S_i^W} p_j p_k
=
\frac{1}{2} \left( L_i^2 - \sum_{j \in S_i^W} p_j^2 \right).
\]

By Cauchy-Schwarz,
\[
L_i^2
=
\left( \sum_{j \in S_i^W} p_j \right)^2
\le
r_i \sum_{j \in S_i^W} p_j^2
\le
m \sum_{j \in S_i^W} p_j^2.
\]

Thus,
\[
\sum_{\{j,k\} \subseteq S_i^W} p_j p_k
\le
\frac{1}{2}(m-1)\sum_{j \in S_i^W} p_j^2.
\]

Summing over all machines,
\[
\sum_{\{j,k\}: j,k \text{ together in } W} p_j p_k
\le
\frac{m-1}{2} S_2.
\]

Hence,
\[
W_{\alpha,m}
\le
\alpha S_2 + \frac{m-1}{2} S_2.
\]

\section{Lower Bound on the Restricted Optimum}

For any schedule,
\[
C_j(\alpha) = B_j + \alpha p_j \ge \alpha p_j,
\]
so
\[
p_j C_j(\alpha) \ge \alpha p_j^2.
\]

Summing over all jobs,
\[
F_\alpha(S) \ge \alpha S_2.
\]

Thus,
\[
OPT_{\alpha,m}^R \ge \alpha S_2.
\]

\section{Approximation Ratio}

Combining the bounds,
\[
\frac{W_{\alpha,m}}{OPT_{\alpha,m}^R}
\le
\frac{\alpha S_2 + \frac{m-1}{2} S_2}{\alpha S_2}
=
1 + \frac{m-1}{2\alpha}.
\]

\section{Tightness}

Consider $n=m$ jobs with equal processing times $p_j = 1$.

\begin{itemize}
\item Worst schedule: all $m$ jobs placed on one machine
\[
W = \alpha m + \binom{m}{2}.
\]

\item Restricted optimal schedule: each job on a different machine
\[
OPT = \alpha m.
\]
\end{itemize}

Thus,
\[
\frac{W}{OPT}
=
\frac{\alpha m + \frac{m(m-1)}{2}}{\alpha m}
=
1 + \frac{m-1}{2\alpha}.
\]

\section{Conclusion}

\[
\boxed{
\frac{W_{\alpha,m}}{OPT_{\alpha,m}^R}
=
1 + \frac{m-1}{2\alpha}
}
\]

This bound is exact and characterizes the worst-case behavior under the pair-separation restriction.

\end{document}